\phantomsection
\section*{Sources and verification scope}
\addcontentsline{toc}{section}{Sources and verification scope}

\subsection*{The mathematical sources}
The realization discussion draws on the companion works listed below. \emph{Komorebi Realization Calculus} supplies the packet, transfer, and defect language. \emph{Komorebi Space} develops coherent views, protected repair, and the distinctions among existence, uniqueness, and stability. \emph{Proof-Path Theory} supplies the certificate and discovery-field distinction, together with its finite models of preservation, provenance, and review. References to recorded verification describe the source reports, not a new replay of those projects.

The elementary calculations in this essay are proved in the text. They illustrate the method and are not presented as new discoveries. The factorization criterion is a set-theoretic statement. The continuation criterion follows from factorization. The return, memory, and defect identities are linear statements. The memory formula retains the declared initial complementary state, including its extra contribution when that state is nonzero. The correction alternative is relative to the specified correction image. The finite-prefix example uses inclusion maps, and the stability example uses the stated bounded operator on \(\ell^2\). The roots-of-unity identity is coefficientwise over \(\mathbb C[[t]]\). Keeping these settings explicit is part of the argument.

\subsection*{What the accompanying checks establish}
The source package includes exact Python checks for the finite examples, the reconstructed memory evolution, coefficient calculations, small matrix identities, and the declared proof-field illustration. It includes deliberate invalid witnesses and false finite claims that those checks must reject. The mathematical tests use exact integer and rational arithmetic. Separate archive tests exercise the integrity validator on empty, incomplete, altered, and malformed records. They corroborate the displayed instances and the implementation, not a universal claim about mathematical discovery.

The proofs concerning an unbounded tower and an infinite Hilbert space are written arguments. Finite tests cannot replace them. No new Lean compilation of the essay or fresh audit of the wider Komorebi corpus is asserted. The accompanying audit record distinguishes calculation checks, source consultation, typesetting, and rendered-page review.

\subsection*{Historical placement}
P\'olya, Lakatos, and experimental mathematics provide intellectual precedents. The Lean sources support the discussion of formal checking. The comparison of research workflows remains a proposed study, not a completed experiment.

\clearpage
\phantomsection
\addcontentsline{toc}{section}{References}
\renewcommand{\refname}{References}
\begin{thebibliography}{9}
\setlength{\itemsep}{8pt}
\bibitem{rc} Komorebi P. \emph{Komorebi Realization Calculus}. Companion monograph essay. See ``Packets and the structure of transport,'' ``The permitted correction determines the answer,'' and ``The residual of a composite must be transported.''
\bibitem{ks} Komorebi P. \emph{Komorebi Space: Coherent Realizations, Obstructions, and Descent}. Companion mathematical monograph. See the preface, the chapters on views and protected correction, and the treatment of approximate answers and their limits.
\bibitem{ppt} Komorebi P. \emph{Proof-Path Theory: Certificates, Compression, Discovery Fields, and Mathematical Memory}. Companion mathematical paper. See Sections 2--4, 7--10, 15--17, and 24.6--24.7 for the specific formal record discussed here.
\bibitem{polya} G. P\'olya. \emph{How to Solve It: A New Aspect of Mathematical Method}. Princeton University Press, 1945. Second edition, 1957. See the prefaces and the discussion of the four phases, pp.~5--6. \href{https://sites.math.rutgers.edu/~zeilberg/math436/polya.pdf}{University-hosted scan}.
\bibitem{lakatos} Imre Lakatos. \emph{Proofs and Refutations: The Logic of Mathematical Discovery}. Edited by John Worrall and Elie Zahar. Cambridge University Press, 1976. \href{https://www.cambridge.org/core/books/proofs-and-refutations/575FC8A6B4FAB79E649EDF5FBB9C6E10}{doi:10.1017/CBO9781139171472}.
\bibitem{expmath} David H. Bailey and Jonathan M. Borwein. ``Experimental Mathematics: Examples, Methods and Implications.'' Author manuscript, 2005. See the introduction, pp.~1--2. \href{https://www.davidhbailey.com/dhbpapers/ams-expmath.pdf}{Author-hosted text}.
\bibitem{leanvalidation} Lean Project. \emph{The Lean Language Reference}, ``Validating a Lean Proof.'' \href{https://lean-lang.org/doc/reference/latest/ValidatingProofs/}{Official documentation}.
\bibitem{leanaxioms} Lean Project. \emph{The Lean Language Reference}, ``Axioms.'' \href{https://lean-lang.org/doc/reference/latest/Axioms/}{Official documentation}.
\bibitem{leancommunity} Lean Community. ``Did you prove it?'' Guidance on project builds, axiom dependencies, and theorem correspondence. \href{https://leanprover-community.github.io/did_you_prove_it.html}{Community verification guidance}.
\end{thebibliography}
